% =====================================================================
%  DRAFT — AI-ASSISTED SCAFFOLD, FOR GROUP REVIEW ONLY.
%  Chapter 1 Contributions entry + Chapter 2 additions (2.7–2.10) + fixes + .bib.
%  NOT for verbatim submission: re-voice in your own words, verify every source
%  against the version of record (K7), and keep the honest wording (no
%  conformance / superiority / novelty claims). Citations use \cite{key} — the
%  matching BibTeX is at the bottom; your bibliographystyle formats them.
% =====================================================================


% #####################################################################
% 1.6  CONTRIBUTIONS  — paste your entry into the (currently empty) 1.6
% #####################################################################
My contribution to the project was the public-facing and staff-facing application layer and its
supporting assurance. I designed and implemented the public discovery application, an accessible,
no-JavaScript-core web interface that renders each result as a typed outcome---supported match,
bounded non-match, or evidence-indeterminate---so that missing or conflicting information is never
presented to the user as a negative fact. I also built the staff assurance application, a role-gated
internal workbench separating the analyst, assurance and authoriser roles, through which staff can
replay a decision's evidence receipts and progress an action through review, approval and send
without being able to skip a step. Underpinning both, I specified the application contracts---the
versioned \texttt{DecisionEnvelope}, the public and staff projections, and the field-level disclosure
matrix---together with the fail-closed projector that enforces them, eight executable
release-blocking invariants, and a fifteen-claim executable assurance case wired to a validator. To
guard the evidence path independently, I implemented a separate certificate checker that verifies
each decision's certificate and witness and rejects forged, stale or over-claimed evidence, tested
against thirteen negative cases and branch mutation. Finally, I carried out the accessibility
evaluation toward WCAG~2.2~AA and prepared the frozen evaluation instruments---the condition manifest
and telemetry schema---used by the evaluation stream.


% #####################################################################
% CHAPTER 2 ADDITIONS  — append after the current §2.6
% #####################################################################

\section{Trust, Reliance and Overreliance in AI-Assisted Decisions}
\label{sec:reliance}
A system that presents machine-derived recommendations to non-expert users must contend not only with
whether its outputs are correct but with how users calibrate their trust in them. The foundational
account of this problem describes trust in automation as an attitude that should be matched to the
automation's actual reliability, so that the design goal is appropriate reliance rather than maximal
reliance \cite{leesee2004trust}. Where reliance is miscalibrated, the classic failure modes are
misuse, the over-reliance on automation that is wrong, and disuse, the rejection of automation that is
right, each carrying its own cost \cite{parasuraman1997humans}. For this project these ideas reframe
the purpose of the interface: it is not to maximise the user's confidence in a recommendation but to
help the user place confidence where the evidence warrants it, and to withhold it where the evidence
does not.

Recent research on human--AI decision-making shows that this calibration does not happen
automatically, and that fluent explanations can make it worse. Empirical work on cognitive-forcing
functions demonstrates that interventions which interrupt automatic acceptance and require the user to
engage can reduce overreliance on AI, albeit sometimes at a cost to user satisfaction
\cite{bucinca2021trust}. Complementary studies show that adding explanations to AI advice does not
reliably improve, and can even harm, complementary team performance, because users may accept
explained-but-incorrect answers more readily \cite{bansal2021whole}. These findings directly motivate
two features of the present design: an explicit confirmation step before a consequential action,
understood as a cognitive-forcing step aimed at appropriate rather than maximal reliance; and an
evidence-first presentation in which the ground for a recommendation is shown alongside it, rather
than a persuasive summary that could invite unwarranted trust.

\section{Explaining Recommendations and Communicating Evidence and Uncertainty}
\label{sec:explanation}
Closely related to reliance is the question of what an explanation in a recommender system is actually
for. A recommendation may be explained in order to increase transparency, to build or calibrate trust,
to improve decision effectiveness, or simply to persuade, and these goals can conflict: an explanation
optimised for persuasion is not the same as one optimised for helping the user judge whether the
recommendation is right for them \cite{tintarev2012evaluating}. This project adopts the
transparency-and-effectiveness reading. The interface's evidence card exists to expose why a candidate
is or is not supported---including a ``why-not'' account for excluded options---so that the
explanation serves the user's judgement rather than the system's persuasiveness.

Where part of the pipeline is generative, explanation also raises the problem of attribution: whether
a natural-language statement is in fact supported by the source it purports to rest on. Work on
measuring attribution in natural-language generation formalises this as whether a system's output can
be traced to, and is entailed by, identified evidence, and provides an evaluation framework for
assessing it \cite{rashkin2023attribution}. This is the standard against which the project's rendering
discipline is set. No factual token about a catalogue attribute is displayed unless it has been
verified against a certified evidence receipt, and states that cannot be attributed are shown as
unknown or conflicting rather than being smoothed into a definite claim. Communicating uncertainty
honestly in this way is treated as a first-class interface requirement rather than as a fallback
behaviour.

\section{Accessibility and Inclusive Interaction}
\label{sec:accessibility}
Because the intended users include people seeking activities under accessibility-related constraints,
the accessibility of the interface is itself a substantive requirement rather than a compliance
afterthought. The Web Content Accessibility Guidelines 2.2 provide the reference success criteria for
perceivable, operable, understandable and robust web content \cite{w3c2023wcag22}. Conformance to
guideline checkpoints is not, however, equivalent to usability for disabled people: empirical
evaluation with blind users has found that a substantial proportion of the problems they actually
encounter are not captured by guideline conformance alone, so that guidelines are ``only half the
story'' \cite{power2012guidelines}. This evidence shapes two decisions in the project. First, the
applications are built on a no-JavaScript core, with routes that remain usable without a map or a
conversational widget, minimising the interaction assumptions placed on assistive technology. Second,
and as a direct consequence, the project deliberately does not claim conformance; it claims to have
been evaluated toward WCAG~2.2~AA within a declared browser and assistive-technology matrix, supported
by manual keyboard and screen-reader testing, precisely because an automated checkpoint pass would
overstate what has been established. A single examiner-found failure would falsify a blanket
conformance claim, whereas the scoped and evidenced claim is defensible.

\section{Positioning: A Capability Audit of Existing Activity-Discovery Tools}
\label{sec:audit}
The design choices above are made against a backdrop of existing tools for discovering physical-activity
opportunities, and the position of this work is best established not by assertion but by a dated
capability audit that compares systems on fixed dimensions. Two broad categories of comparator are
relevant. The first is consumer-facing discovery and booking services built on the same open-data
ecosystem, of which OpenActive-aligned activity finders such as Get Active are representative; the
second is commercial aggregation and booking infrastructure that consolidates provider feeds, of which
imin is an example. Alongside these sit grounded question-answering assistants from the public sector
and data-grounded generation efforts more broadly, which share the reliability concerns discussed in
Section~\ref{sec:grounded}.

The audit compares these systems on dimensions that follow directly from the project's aims rather than
on general feature counts: whether a tool distinguishes an activity being absent from the observed
catalogue from the tool simply lacking the evidence to decide; whether it communicates uncertainty and
provenance to the user or presents all results with uniform confidence; whether any factual claim it
displays is independently checkable; whether it abstains or qualifies rather than answering when
evidence is insufficient; and whether its accessibility is evaluated rather than assumed. Because these
dimensions require inspecting each system as deployed at a specific date, the audit is maintained as
dated evidence rather than reported from memory, and this dissertation does not claim priority, novelty
or superiority over any named incumbent in advance of that comparison. What the audit establishes is
the position of the contribution: the specific combination of evidence-bounded decisions, independently
checkable receipts, and an interface disciplined not to convert absence into denial is the gap the
present design occupies. Establishing that this combination outperforms an incumbent would require the
controlled comparison the evaluation instruments are built to support, and is not claimed here.


% #####################################################################
% FIXES to the EXISTING Chapter 2 text (find-and-replace)
% #####################################################################
% (1) §2.2 — remove the editing marker. REPLACE:
%     "...to support a query result (Razniewski & Nutt, 2014) Need to check properly."
%   WITH:
%     "...to support a given query result, independent of which values happen to be present
%      \cite{razniewski2014completeness}."
%   THEN verify that Razniewski & Nutt reference against the version of record and add its bib entry.
%
% (2) §2.5 — the forward reference now points to the audit section. REPLACE:
%     "examined further in the capability audit (Section 2.9)."
%   WITH:
%     "examined further in the capability audit (Section~\ref{sec:audit})."
%
% (3) Citation style — the chapter mixes [n] and (Author, year). Cite everything with \cite{key}
%   and let one \bibliographystyle format them consistently; do not hand-type "(Author, year)".
%
% (4) §2.10 cross-refs §2.5 via \ref{sec:grounded}. Add a label to your existing §2.5 heading:
%     \section{Grounded Conversational AI and its Reliability Limits}\label{sec:grounded}
%   (or just write "(Section 2.5)" as plain text instead of the \ref).


% #####################################################################
% BIBTEX — add these to your .bib file.  (bucinca / rashkin / power / w3c
% are already in execution-refs.bib; reuse those keys, don't duplicate.)
% DOIs verified against publisher records; the 4 new entries are below.
% #####################################################################
%
% @article{leesee2004trust,
%   author  = {Lee, John D. and See, Katrina A.},
%   title   = {Trust in Automation: Designing for Appropriate Reliance},
%   journal = {Human Factors}, volume = {46}, number = {1}, pages = {50--80},
%   year    = {2004}, doi = {10.1518/hfes.46.1.50_30392} }
%
% @article{parasuraman1997humans,
%   author  = {Parasuraman, Raja and Riley, Victor},
%   title   = {Humans and Automation: Use, Misuse, Disuse, Abuse},
%   journal = {Human Factors}, volume = {39}, number = {2}, pages = {230--253},
%   year    = {1997}, doi = {10.1518/001872097778543886} }
%
% @inproceedings{bansal2021whole,
%   author    = {Bansal, Gagan and Wu, Tongshuang and Zhou, Joyce and Fok, Raymond and Nushi, Besmira and Kamar, Ece and Ribeiro, Marco Tulio and Weld, Daniel S.},
%   title     = {Does the Whole Exceed its Parts? The Effect of {AI} Explanations on Complementary Team Performance},
%   booktitle = {Proceedings of the 2021 CHI Conference on Human Factors in Computing Systems},
%   year      = {2021}, doi = {10.1145/3411764.3445717} }
%
% @article{tintarev2012evaluating,
%   author  = {Tintarev, Nava and Masthoff, Judith},
%   title   = {Evaluating the Effectiveness of Explanations for Recommender Systems},
%   journal = {User Modeling and User-Adapted Interaction}, volume = {22}, number = {4--5},
%   pages   = {399--439}, year = {2012}, doi = {10.1007/s11257-011-9117-5} }
